\section{Worked examples}
\label{sec:illustrations}

Each example is a diagram on which \cref{thm:lp} is a linear program and no bipartition meets Proposition~2 of \citet{symbolicbounds2023general}. The failure of each bipartition is recorded in \cref{app:sachs}. The query is a counterfactual event or a linear combination of such events. The interval returned by \texttt{bound} is the identified range.

\subsection{A pinned mediator}
\label{sec:ill-mediated}
\label{sec:ill-solver}

On $\G:\{X\to W\to Y,\, X\leftrightarrow W\}$ the latent of $\{X,W\}$ is independent of the noise of $Y$. The table is
\begin{equation}
\label{eq:re3}
P(x,w)=\begin{pmatrix}0.25 & 0.25\\ 0.15 & 0.35\end{pmatrix},
\qquad
P(Y{=}1\mid W{=}w)=(0.4,\,0.6),
\end{equation}
with rows indexed by $x$ and $Y\perp X\mid W$. The query is the combination
\[
P(Y_{X=1}{=}1)-P(Y_{X=0}{=}1).
\]
Each event evaluates $Y$ in one world, so $Y$ is pinned to $P(Y\mid W)$ in both, the remaining component $\{X,W\}$ has one hyperedge, and the difference is linear in that one vector. The sharp interval is $[-0.080,0.120]$. The single term $P(Y_{X=1}{=}1)$ is the LP of \cref{ex:mediated}, with interval $[0.470,0.570]$.

$Y$ does not share the latent of $\{X,W\}$, so no bipartition meets their conditions. Replacing the latent by one hyperedge on $\{X,W,Y\}$ drops the independence and widens the combination to $[-0.480,0.520]$. The cross-world event $P(Y_{X=1}{=}1,\,Y_{X=0}{=}0)$ evaluates $Y$ in two worlds, so $Y$ is not pinned and the program is no longer linear.

\subsection{An instrument confounded with the treatment}
\label{sec:ill-designs}

A randomized offer $Z$ encourages enrollment $X$, with outcome $Y$. The law of $(Z,X,Y)$, in that order, is
\begin{equation}
\label{eq:iv-table}
\begin{aligned}
P(0,0,0)&=0.3543,& P(0,0,1)&=0.1173,& P(0,1,0)&=0.2024,& P(0,1,1)&=0.0860,\\
P(1,0,0)&=0.0377,& P(1,0,1)&=0.0421,& P(1,1,0)&=0.0923,& P(1,1,1)&=0.0679.
\end{aligned}
\end{equation}
The query is the joint counterfactual $P(Y_{X=1}{=}1,\,Z{=}0)$. On $\G:\{Z\to X\to Y,\, X\leftrightarrow Y\}$ the instrument lies on no hyperedge and is evaluated once, so it is pinned, and the interval is $[0.215,0.468]$, namely $P(Z{=}0)$ times the Balke--Pearl range $[0.283,0.615]$ for $P(Y_{X=1}{=}1)$; this graph meets their Proposition~2 for the marginal.

On $\G:\{Z\to X\to Y,\, Z\leftrightarrow X,\, X\leftrightarrow Y\}$ the offer shares a latent with treatment and not with the outcome. The variable $Z$ is a root of that component, so the independence of the two latents is the linear constraint of \cref{prop:lift}. Both intervals $[0.215,0.468]$ and $[0.283,0.615]$ persist (whereas one hyperedge on $\{Z,X,Y\}$ drops that independence and widens the joint to $[0.086,0.558]$). An anchored odds-ratio bound of $\Lambda=3$ on the propensity of $X$ adds linear inequalities in that response law whose coefficients are fixed by the table \citep{tan2006distributional}. The program remains linear (\cref{prop:assumptions}), and the joint interval narrows to $[0.251,0.370]$. Every bipartition is met by the latent on $\{Z,X\}$ or by the latent on $\{X,Y\}$, so Proposition~2 does not apply.

\subsection{Two instruments}
\label{sec:ill-two}

On $\G:\{Z_1\to X,\, Z_2\to X,\, X\to Y,\, X\leftrightarrow Y\}$ the instruments are independent. The law of $(Z_1,Z_2,X,Y)$ is
\begin{equation}
\label{eq:two-iv}
\begin{aligned}
P(0,0,0,0)&=0.0159,& P(0,0,0,1)&=0.0062,& P(0,0,1,0)&=0.0059,& P(0,0,1,1)&=0.0069,\\
P(0,1,0,0)&=0.0254,& P(0,1,0,1)&=0.0084,& P(0,1,1,0)&=0.0122,& P(0,1,1,1)&=0.0117,\\
P(1,0,0,0)&=0.0953,& P(1,0,0,1)&=0.0439,& P(1,0,1,0)&=0.0669,& P(1,0,1,1)&=0.1358,\\
P(1,1,0,0)&=0.2643,& P(1,1,0,1)&=0.1175,& P(1,1,1,0)&=0.0711,& P(1,1,1,1)&=0.1126.
\end{aligned}
\end{equation}
The query is the combination
\[
P(Y_{X=1}{=}1,\,Z_1{=}1)+P(Y_{X=0}{=}0,\,Z_2{=}1).
\]
Each event evaluates one instrument in the factual world and the outcome in one world. Both instruments lie on no hyperedge, so each is pinned, and the objective is linear in the response law of $\{X,Y\}$. The sharp interval is $[0.652,1.210]$ (versus $[0.538,1.267]$ under one hyperedge on $\{Z_1,Z_2,X,Y\}$).

The cut with all arrows pointing the same way puts both instruments on the left and $\{X,Y\}$ on the right. Their condition that the left side share one latent fails, since $Z_1\perp Z_2$.

\subsection{A pair chain}
\label{sec:ill-pair}

On $\G:\{X\to W\to Y,\, X\leftrightarrow W,\, W\leftrightarrow Y\}$ the joint law of $(X,W,Y)$ is
\begin{equation}
\label{eq:re2}
\begin{aligned}
P(0,0,0)&=0.20,& P(0,0,1)&=0.05,& P(0,1,0)&=0.10,& P(0,1,1)&=0.15,\\
P(1,0,0)&=0.05,& P(1,0,1)&=0.10,& P(1,1,0)&=0.10,& P(1,1,1)&=0.25,
\end{aligned}
\end{equation}
whose $(X,Y)$-margin is
\begin{equation}
\label{eq:re1}
\begin{aligned}
P(X{=}0,Y{=}0)&=0.30,& P(X{=}0,Y{=}1)&=0.20,\\
P(X{=}1,Y{=}0)&=0.15,& P(X{=}1,Y{=}1)&=0.35.
\end{aligned}
\end{equation}
The variable $X$ is a root, so covering holds and \cref{prop:lift} makes every linear query an LP in the joint response types. The combination $P(Y_{X=1}{=}1)+P(W_{X=1}{=}1)$ has sharp interval $[0.800,1.600]$. The cross-world event $P(Y_{X=1}{=}1,\,Y_{X=0}{=}0)$ has sharp interval $[0.000,0.600]$.

One hyperedge on $\{X,W,Y\}$ meets their Proposition~2. It widens the combination to $[0.750,1.600]$ and the cross-world event to $[0.000,0.650]$. The two-latent graph has no such bipartition: neither latent contains all three variables, and each latent meets both sides of the cut that would separate its endpoints (\cref{ex:pair}).

Monotone response of $W$ in $X$ sets to zero every response of $W$ that decreases in $X$ \citep{manski1997monotone}. The restriction is a face of the remaining response law, so the program stays linear (\cref{prop:assumptions}). The combination narrows to $[1.150,1.600]$ and the cross-world event to $[0.000,0.450]$.
